\input{../tex-inputs/latex-leseansicht-vorspann}

\section[Gustav Schwarzkopf an Arthur Schnitzler, 12. 7. 1912]{L04264 Gustav Schwarzkopf an Arthur Schnitzler, 12. 7. 1912}
\nopagebreak\mylabel{L04264v}
\rehead{ }\normalsize\beginnumbering\briefempfaengerindex{Schnitzler, Arthur@\textsc{Schnitzler, Arthur}!zzzSchwarzkopf, Gustav@\emph{von Gustav Schwarzkopf}!1912-07-121@{12. 7. 1912}|(be}
\toendnotes[C]{\smallbreak\pagebreak[2]}
\correspDesc{Versand  durch Gustav Schwarzkopf am 12. 7. 1912 in Reifnitz [Maria Wörth]
\newline{}Erhalt  durch Arthur Schnitzler im Zeitraum [13. 7. 1912 – 17. 7. 1912?] in Wien}\toendnotes[C]{\smallbreak}
\Standort{CUL, Schnitzler, B 96a.}
\physDesc{Bildpostkarte, 384 Zeichen
\newline{}Handschrift: Bleistift, deutsche Kurrent
\newline{}Versand: 1) Stempel: »\nobreak{}\oindex{Reifnitz [Maria Wörth]@\textbf{Reifnitz [Maria Wörth]}|pwk}Reifnitz am Wörthersee, 13. VII. 12\nobreak{}«.   2) mit Bleistift von unbekannter Hand Streichung der falschen Bezirksangabe in der Adressierung: »XIX«}\toendnotes[C]{\smallbreak}\pstart{}{\pb}Herrn\pend{}\pstart{}\textsc{D\textsuperscript{r} Arthur Schnitzler}\pend{}\pstart{}\textsc{Wien XIX\oindex{XIX., Döbling@\textbf{XIX., Döbling}, \emph{Verwaltungsgebiet}|pw}}\pend{}\pstart{}Sternwarteſtraße 71\oindex{Wien@\textbf{Wien}!XVIII., Währing@\textbf{XVIII., Währing}!Sternwartestraße 71@\textbf{Sternwartestraße 71}, \emph{Wohngebäude}|pw}\pend{}{\bigskip}
\pstart
           \noindent{}\centering{}{\pb}\textcolor{gray}{\textbf{Reifnitz am Wörthersee\oindex{Reifnitz [Maria Wörth]@\textbf{Reifnitz [Maria Wörth]}|pw}.}}\pend
           \vspace{1em}
\pstart
           \raggedleft{}Reifnitz\oindex{Reifnitz [Maria Wörth]@\textbf{Reifnitz [Maria Wörth]}|pw}{ }12. VII.\pend
           
\pstart{}Lieber Arthur!\pend\vspace{0.5em}
\pstart
           \textcolor{gray}{N}un bin ich hier geſtrandet. Erſatz für
              Brioni\oindex{Brijuni@\textbf{Brijuni}|pw}. Der berühmte
              \label{K_L04264-1v}\edtext{Wiener See\oindex{Wörthersee@\textbf{Wörthersee}, \emph{See}|pwv}}{\lemma{\textnormal{\emph{Wiener See}}}\Cendnote{\textnormal{unklare Anspielung}}}\label{K_L04264-1} hat aber nur
      14 ½°. Für Abwechslung iſt
      durch tägliche Gewitter geſorgt.
      Aber dafür wohne ich
      in der »\textsc{Villa Grande\oindex{Villa Grande@\textbf{Villa Grande}, \emph{Beherbergungsgebäude}|pw}}«
      und hat ein Zi{\geminationm}er mit
      \uline{Balkon}! Sehr nett. Nur
      umdrehen darf man ſich nicht.\pend
           
\pstart
           Ihre Frau\pwindex{Schnitzler, Olga 17.\,1.\,1882 Wien – 13.\,1.\,1970 Lugano@\textsc{Schnitzler, Olga} (17.\,1.\,1882 Wien – 13.\,1.\,1970 Lugano), \emph{Schauspielerin, Sängerin}|pwv} und Sie herzlichſt{\\[\baselineskip]}grüßend{\\[\baselineskip]}\spacefill\mbox{Gustav Sch.}\pend
           \leftskip=0em{}\selectlanguage{ngerman}\endnumbering\briefempfaengerindex{Schnitzler, Arthur@\textsc{Schnitzler, Arthur}!zzzSchwarzkopf, Gustav@\emph{von Gustav Schwarzkopf}!1912-07-121@{12. 7. 1912}|)be}\mylabel{L04264h}
\begin{anhang}
\end{anhang}\newcommand{\dateiname}{L04264}\newcommand{\titel}{Gustav Schwarzkopf an Arthur Schnitzler, 12. 7. 1912}\newcommand{\editorInnen}{Herausgegeben von Jahnke, SelmaMüller, Martin Anton}\input{../tex-inputs/latex-leseansicht-abspann}
